\chapter{Affidabilit\`a e scalabilit\`a}
\label{chap:reliability}

\section{Meccanismi di affidabilit\`a adottati}
\label{sec:reliability_meccanismi}

I meccanismi di affidabilit\`a sono stati selezionati in modo deliberato e motivato,
scegliendo un sottoinsieme coerente con il profilo di carico e il dominio del servizio.

\subsection{Idempotenza forte sull'ingestione email}

Il vincolo \texttt{UNIQUE(email\_account\_id, message\_uid)} nella tabella
\texttt{email\_jobs} impedisce la creazione di job duplicati anche in presenza di:
\begin{itemize}[leftmargin=*]
    \item Riavvii del poller a met\`a ciclo.
    \item Esecuzioni concorrenti di pi\`u istanze del poller.
    \item Network timeout che causano retry del messaggio gi\`a inserito.
\end{itemize}

Il campo \texttt{last\_uid} per account funge da \emph{watermark monotono}: viene
aggiornato al massimo UID osservato anche nei rami di errore. Questo garantisce la
propriet\`a \emph{at-least-once} (ogni messaggio viene processato almeno una volta)
senza degenerare in \emph{at-most-once} (nessun messaggio viene perso in caso di
restart).

\subsection{Initial-skip sulle nuove cassette}

Al primo poll di un account (\texttt{last\_synced\_at IS NULL}), il poller imposta
\texttt{last\_uid} al massimo UID corrente, evitando di scaricare l'intera storia della
casella email. Questo contiene il consumo iniziale di banda, di chiamate Tika e di
token LLM per account con anni di storia.

\subsection{Connection pooling PostgreSQL}

Il motore SQLAlchemy \`e configurato con:

\begin{lstlisting}[language=Python, caption={Connection pooling (\texttt{backend/app/database.py}).}]
engine = create_engine(
    db_url,
    poolclass=QueuePool,
    pool_size=20,        # connessioni aperte di base
    max_overflow=10,     # connessioni aggiuntive in picco
    pool_recycle=3600,   # riciclo orario per evitare connessioni stantie
    pool_pre_ping=True,  # verifica la connessione prima di usarla
    echo=False
)
\end{lstlisting}

Il parametro \texttt{pool\_pre\_ping=True} emette un \texttt{SELECT 1} prima di ogni
utilizzo di una connessione dal pool, garantendo la rimozione automatica di connessioni
morte (classe di errori ``server closed the connection unexpectedly'').

\subsection{Pattern two-phase upload (presigned URL)}

Il dato binario non transita dal backend: il backend agisce esclusivamente come piano
di controllo (genera il presigned URL) mentre il client carica direttamente su SeaweedFS.
Questo pattern:
\begin{itemize}[leftmargin=*]
    \item Riduce il traffico verso il backend, che non \`e un collo di bottiglia per
          gli upload anche sotto carico elevato.
    \item Abilita lo scaling indipendente del backend rispetto al volume di dati.
    \item Previene la classe di errori ``record orfano'' tramite la fase di confirm.
\end{itemize}

\subsection{Streaming download}

I download di file avvengono tramite \texttt{StreamingResponse} con \texttt{iter\_chunks}:

\begin{lstlisting}[language=Python, caption={Download in streaming (\texttt{backend/app/routes/documents.py}).}]
return StreamingResponse(
    file_data,
    media_type=doc.mime_type or "application/octet-stream",
    headers={"Content-Disposition": f'attachment; filename="{doc.name}"'}
)
\end{lstlisting}

Questo evita di caricare in memoria l'intero contenuto del file prima della risposta,
mantenendo costante la pressione di memoria del backend anche per file di grandi dimensioni.

\subsection{Graceful degradation dei worker}

La funzione \texttt{\_process\_account} nel poller aggiorna \texttt{last\_synced\_at}
anche in caso di errore, evitando un \emph{retry storm} a 60~s su un account
permanentemente in down (credenziali revocate, IMAP irraggiungibile). L'errore viene
loggato ma non interrompe il ciclo sugli altri account.

Nell'agent worker, le funzioni \texttt{\_process\_job\_safe} e
\texttt{\_process\_manual\_upload\_safe} incapsulano la logica di elaborazione con
gestione dell'eccezione:

\begin{lstlisting}[language=Python, caption={Error handling nel worker (\texttt{agent\_worker/worker.py}).}]
def _process_job_safe(db, job, os_client):
    try:
        job.status = "processing"
        db.commit()
        _process_job(db, job, os_client)
    except Exception as e:
        db.rollback()
        logger.error("Job %s failed: %s", job.id, e, exc_info=True)
        try:
            job.status = "failed"
            job.error_message = str(e)
            job.processed_at = datetime.now(timezone.utc)
            db.commit()
        except Exception:
            db.rollback()
\end{lstlisting}

Il loop principale del worker prosegue indipendentemente dal fallimento del singolo job.

\subsection{Transactional safety}

Tutte le scritture che toccano pi\`u tabelle sono racchiuse in transazioni esplicite.
Il pattern \texttt{db.flush()} viene usato per ottenere UUID prima del commit finale,
evitando l'utilizzo di chiavi non ancora persistite:

\begin{lstlisting}[language=Python, caption={Uso di flush per ottenere UUID pre-commit.}]
doc = Document(name=att.filename, ...)
db.add(doc)
db.flush()   # ottieni doc.id senza commit

agent_file = AgentFile(document_id=doc.id, ...)
db.add(agent_file)
db.flush()   # ottieni agent_file.id

# ... operazioni S3 ...
db.commit()  # commit atomico di tutto
\end{lstlisting}

\subsection{Automatic bucket provisioning}

Il modulo \texttt{storage.py} verifica l'esistenza del bucket prima di ogni operazione
e lo crea in modalit\`a \emph{lazy} se assente:

\begin{lstlisting}[language=Python]
def ensure_bucket(s3_client, bucket_name: str) -> None:
    try:
        s3_client.head_bucket(Bucket=bucket_name)
    except ClientError:
        s3_client.create_bucket(Bucket=bucket_name)
\end{lstlisting}

Questo elimina la classe di errori operativi al primo avvio su un ambiente nuovo.

\subsection{Deduplicazione contenuti}

Prima di classificare un allegato, l'agent worker calcola l'hash SHA-256 del contenuto
e lo salva su \texttt{agent\_files.content\_hash}. File identici ricevuti pi\`u volte
(ad es.\ la stessa fattura inoltrata) non vengono ri-classificati, con risparmio di
token LLM e prevenzione di archiviazioni duplicate.

\subsection{Paginazione}

Tutti gli endpoint di listing espongono i parametri \texttt{limit} (1--100) e \texttt{offset},
prevenendo query senza bounds che potrebbero saturare la memoria del backend su tenant
con migliaia di documenti.

\section{Meccanismi di scalabilit\`a}
\label{sec:scalabilita}

\subsection{Scaling orizzontale del backend}

Il backend FastAPI \`e \emph{stateless} per costruzione: non mantiene stato di sessione
in memoria. L'unica eccezione \`e la cache \texttt{\_last\_updated} del
\texttt{LastActiveMiddleware}, accettata perch\'e il debounce di 5~minuti \`e tollerante
a deriva tra istanze diverse. Il backend pu\`o essere replicato orizzontalmente dietro
un load balancer senza modifiche al codice.

\subsection{Scaling dell'agent worker}

Il design della macchina a stati su \texttt{email\_jobs} (\texttt{pending} $\rightarrow$
\texttt{processing}) \`e concorrenza-safe: due istanze del worker che prelevano lo stesso
job in modo concorrente produrranno al pi\`u un duplicato transazionale, risolto dal
vincolo UNIQUE sull'output. Pi\`u istanze del worker possono operare in parallelo su
job diversi.

\subsection{Scaling dell'object storage}

SeaweedFS scala aggiungendo volume server al cluster esistente. Il sharding dei dati
avviene a livello di file (ogni file \`e memorizzato su un volume server in base a
regole di bilanciamento del master). Aggiungere capacit\`a non richiede modifiche
all'applicazione, che interagisce sempre tramite l'interfaccia S3 standard.

\subsection{Scaling del database}

Una singola istanza PostgreSQL \`e sufficiente per il carico del MVP. La separazione
tra write-heavy (worker e poller che inseriscono job e operazioni) e read-heavy
(browse UI, ricerche) consente in futuro l'introduzione di una \emph{read replica}
senza modifiche applicative significative, indirizzando le query di lettura alla replica.

\subsection{Scaling della ricerca}

Gli indici OpenSearch sono per-tenant: ogni utente ha il proprio indice indipendente.
Questo consente di bilanciare il carico di ricerca tra nodi OpenSearch e di eliminare
un indice di tenant senza impattare gli altri. \TODO{Configurazione del cluster OpenSearch
multi-nodo per alta disponibilit\`a}
